% Options for packages loaded elsewhere
\PassOptionsToPackage{unicode}{hyperref}
\PassOptionsToPackage{hyphens}{url}
%
\documentclass[
  brazilian,
]{article}
\usepackage{lmodern}
\usepackage{amssymb,amsmath}
\usepackage{ifxetex,ifluatex}
\ifnum 0\ifxetex 1\fi\ifluatex 1\fi=0 % if pdftex
  \usepackage[T1]{fontenc}
  \usepackage[utf8]{inputenc}
  \usepackage{textcomp} % provide euro and other symbols
\else % if luatex or xetex
  \usepackage{unicode-math}
  \defaultfontfeatures{Scale=MatchLowercase}
  \defaultfontfeatures[\rmfamily]{Ligatures=TeX,Scale=1}
\fi
% Use upquote if available, for straight quotes in verbatim environments
\IfFileExists{upquote.sty}{\usepackage{upquote}}{}
\IfFileExists{microtype.sty}{% use microtype if available
  \usepackage[]{microtype}
  \UseMicrotypeSet[protrusion]{basicmath} % disable protrusion for tt fonts
}{}
\makeatletter
\@ifundefined{KOMAClassName}{% if non-KOMA class
  \IfFileExists{parskip.sty}{%
    \usepackage{parskip}
  }{% else
    \setlength{\parindent}{0pt}
    \setlength{\parskip}{6pt plus 2pt minus 1pt}}
}{% if KOMA class
  \KOMAoptions{parskip=half}}
\makeatother
\usepackage{xcolor}
\IfFileExists{xurl.sty}{\usepackage{xurl}}{} % add URL line breaks if available
\IfFileExists{bookmark.sty}{\usepackage{bookmark}}{\usepackage{hyperref}}
\hypersetup{
  pdftitle={Parte III: Processamento de Linguagem Natural para Linguistas},
  pdfauthor={CiberExt 26-29 · FEELT38103 · Universidade Federal de Uberlândia},
  pdflang={pt-BR},
  hidelinks,
  pdfcreator={LaTeX via pandoc}}
\urlstyle{same} % disable monospaced font for URLs
\setlength{\emergencystretch}{3em} % prevent overfull lines
\providecommand{\tightlist}{%
  \setlength{\itemsep}{0pt}\setlength{\parskip}{0pt}}
\setcounter{secnumdepth}{-\maxdimen} % remove section numbering
\ifxetex
  % Load polyglossia as late as possible: uses bidi with RTL langages (e.g. Hebrew, Arabic)
  \usepackage{polyglossia}
  \setmainlanguage[]{brazilian}
\else
  \usepackage[shorthands=off,main=brazilian]{babel}
\fi

\title{Parte III: Processamento de Linguagem Natural para Linguistas}
\usepackage{etoolbox}
\makeatletter
\providecommand{\subtitle}[1]{% add subtitle to \maketitle
  \apptocmd{\@title}{\par {\large #1 \par}}{}{}
}
\makeatother
\subtitle{Unidade 2 - Respostas dos Exercícios}
\author{CiberExt 26-29 · FEELT38103 · Universidade Federal de
Uberlândia}
\date{Agosto de 2026}

\begin{document}
\maketitle

\hypertarget{gabarito---muxf3dulo-2-dependuxeancias-universais}{%
\subsection{Gabarito - Módulo 2: Dependências
Universais}\label{gabarito---muxf3dulo-2-dependuxeancias-universais}}

\textbf{Resposta 1:} 1. Desenvolver um framework comum para descrever a
estrutura gramatical de diversos idiomas. 2. Criar corpora anotados (ou
\emph{treebanks}) para diversos idiomas aplicando este framework
universal.

\textbf{Resposta 2:} 1. \textbf{Nominais (\emph{nominals}):} Usados para
representar coisas, sendo frequentemente construídos em torno de
substantivos (análogos a sintagmas nominais). 2. \textbf{Orações
(\emph{clauses}):} Usadas para representar eventos, sendo construídas
fundamentalmente em torno de verbos. Elas situam-se acima dos nominais
hierarquicamente. 3. \textbf{Modificadores (\emph{modifiers}):} Dependem
de adjetivos e advérbios e servem para expandir, refinar e detalhar o
significado tanto de nominais quanto de orações.

\textbf{Resposta 3:} A etiqueta \texttt{root} identifica a \textbf{raiz}
ou a cabeça central da oração (geralmente o verbo principal). No
framework UD, as dependências são desenhadas de cabeça para cabeça;
portanto, os arcos sintáticos partem da raiz (o verbo) em direção às
palavras que atuam como as cabeças dos nominais dependentes (como o
pronome sujeito ou o substantivo objeto direto).

\textbf{Resposta 4:} O atributo \texttt{token.subtree} retorna um
gerador que contém todos os dependentes diretos e indiretos (filhos,
netos, bisnetos, etc.) \textbf{incluindo o próprio \emph{Token}
original}. O atributo \texttt{token.children}, por outro lado, retorna
apenas os dependentes imediatos (filhos diretos de primeiro nível) e
\textbf{não inclui o \emph{Token} original} na sua saída.

\textbf{Resposta 5:} Essas métricas penalizam desproporcionalmente
idiomas morfologicamente ricos. Em inglês, as sentenças são mais longas
devido ao alto uso de palavras de função independentes (preposições,
auxiliares); logo, errar uma palavra causa um pequeno impacto
percentual. Em finlandês, onde essas funções são acopladas nas palavras
originais (formando sentenças mais curtas), um único erro sintático
representa uma porcentagem de perda muito maior na métrica final.
Métricas alternativas sugeridas incluem a \textbf{CLAS}
(\emph{content-labeled attachment score}), que ignora palavras de função
na contagem, a \textbf{MLAS} (com foco na morfologia) e a \textbf{BLEX}
(com foco no lema).

\begin{center}\rule{0.5\linewidth}{0.5pt}\end{center}

\end{document}
